\section{The first three quantum counterterms}
\label{sec:first-counterterms}

Throughout this section $\g=\fsl_N$ and $\kappa=\tr$.
To retain the original cochain notation, write
\[
 D(X_0e^{p_0z},X_1e^{p_1z},X_2e^{p_2z})
   =[p_0p_1]\tr(X_0\{X_1,X_2\}),\qquad
 p_0+p_1+p_2=0.
\]
Here $D$ denotes a cubic cochain, whereas the $D$ acting on
superfields in \eqref{eq:classical} is the free differential.
The matrix anticommutator is
$\{X_1,X_2\}=X_1X_2+X_2X_1$.
The color tensor is symmetric and the momentum determinant is
alternating, so the complete cochain is alternating.

\subsection{The forced order-one interaction}

At first order the only invariant equivariant cubic is $j=aD$.
Its classical differential vanishes. Contraction of the full
shuffle bracket, using the traceless completeness relation, gives
\begin{equation}
 \frac12\{D,D\}
  =\frac{F_{223}}{12N^2}
     +\frac{N^2+4}{2N^2}\,\beta .
 \label{eq:DD}
\end{equation}
Substitution in \eqref{eq:first-two-equations} yields
\begin{align}
 d i_2
 &=\left(\frac{a^2}{12N^2}
                       +\frac{\epsilon\pi\nu}{48}\right)F_{223}
   +\left(\frac{a^2(N^2+4)}{2N^2}
                       -\frac{\epsilon\pi\nu N^2}{40}\right)\beta .
 \label{eq:order-two-assembly}
\end{align}
For $N\ge3$, Theorem~\ref{thm:bare-anomaly} forces the first
coefficient to vanish. Therefore
\begin{align}
 a^2&=-\frac{\epsilon\pi\nu N^2}{4}
       =\frac{\epsilon U N^2}{256\pi^4},\notag\\
 i_2&=\frac{\epsilon U(3N^2+10)}{1280\pi^4}
         [p_0p_1]^2\tr(X_0[X_1,X_2]).
 \label{eq:first-two-counterterms}
\end{align}
These are equalities of full cochains, not only cohomology classes.
In the convention $\epsilon U=1$, the two roots are
$a=N/(16\pi^2)$ and $a=-N/(16\pi^2)$.

\subsection{The marked two-loop contribution at order three}

There are three sources at order $\hbar^3$: the classical
differential of $i_3$, the tree bracket of $j$ and $i_2$, and the
two-loop theta graphs with one $j$ vertex. Bare three-loop graphs
and one-loop graphs with any of these vertices vanish.
The complete differentiated theta calculation gives
\begin{align}
 \{D,[p_0p_1]^2f\}
  &=\frac13d\bigl([p_0p_1]^3\tr(X_0\{X_1,X_2\})\bigr),\notag\\
 \Theta(D)
  &=-\frac{N^2}{60}
       d\bigl([p_0p_1]^3\tr(X_0\{X_1,X_2\})\bigr).
 \label{eq:marked-theta}
\end{align}
Here $\Theta(D)$ is the sum over both theta graphs and all
single marked vertices, with the common scalar $\pi\nu$
removed. Its graph automorphism factors are included.
The local equation is consequently
\[
 d i_3=\{j,i_2\}+\epsilon\pi\nu\,a\,\Theta(D).
\]
Using \eqref{eq:first-two-counterterms} and \eqref{eq:marked-theta},
\begin{align}
 i_3
 &=\frac{a\epsilon U}{\pi^4}
       \left(\frac{3N^2+10}{3840}+\frac{N^2}{3840}\right)
          [p_0p_1]^3\tr(X_0\{X_1,X_2\})\notag\\
 &=\frac{a\epsilon U(2N^2+5)}{1920\pi^4}
          [p_0p_1]^3\tr(X_0\{X_1,X_2\}).
 \label{eq:i3}
\end{align}
The tensor identities in this section were independently replayed
from the original graphs using exact rational matrices; their
finite evaluations are certificates accompanying the formal
trace identities, rather than substitutes for them.

\subsection{Rank two and the two branches}

For traceless $2\times2$ matrices,
\[
 X_1X_2+X_2X_1=\tr(X_1X_2)\mathbf1_2,\qquad
 \tr(X_0\{X_1,X_2\})=0 .
\]
Evaluation of the full identities at $N=2$ therefore gives
\begin{equation}
 j=0,\qquad
 i_2=\frac{11\epsilon U}{640\pi^4}
             [p_0p_1]^2\tr(X_0[X_1,X_2]),\qquad
 i_3=0.
 \label{eq:rank-two}
\end{equation}
It is not necessary to cancel a nonzero class at rank two:
the bare anomaly there is already exact.

The Lie automorphism $X\mapsto-X^{\mathsf T}$ preserves the trace
pairing and fixes $f$, while reversing the symmetric cubic.
It commutes with the cochain differential and every graph
contraction. Once the all-order uniqueness in the formal sector
is proved, it follows that the two choices of $a$ are exchanged
by this automorphism: every even-order counterterm is fixed and
every odd-order counterterm changes sign. Thus for $N\ge3$
there are two leading couplings in the fixed color convention,
related by a symmetry of the underlying classical theory.
At $N=2$ every odd-order cubic vanishes, so the two formal
branches specialize to the same quantization.
